\chapter*{Preface}

An algorithm is a proof that does something. To claim that a method sorts a list, finds the shortest route through a
network or tells a prime from a composite is to state a theorem, and to claim that it does so quickly is to state a
second one. This book proves both theorems, for the algorithms that every program relies on and for some of the most
beautiful ideas computer science has produced.

We begin by learning to measure: recurrences, amortised costs, and the lower bounds that say no cleverness can do
better. Then come sorting and the networks that sort without making decisions, the data structures that make
everything else fast (heaps, balanced trees, hashing, union--find), the classical algorithms on graphs, and dynamic
programming. The middle of the book turns the existence theorems of Tome~C1 into efficient procedures: flows,
matchings in bipartite and general graphs, and linear programming with its duality, which explains why so many
optimisation problems come with a certificate of optimality. Randomness enters next, as a tool: random contractions
find minimum cuts, random evaluations test polynomial identities, a few random bits summarise an enormous stream.
The last chapters are algebraic and combinatorial: fast multiplication of numbers, polynomials and matrices, primality
testing up to the deterministic algorithm of Agrawal, Kayal and Saxena, the algorithms that search texts and
genomes, and, when exact answers are out of reach, approximation and online algorithms with proved guarantees.

The prerequisites are Tomes~1, 2 and~3 and Tome~C1; results from them are recalled where they are used, with a
precise reference. Every theorem is proved, and every running time is a theorem. Each chapter opens with a question
to think about, and ``Pause and think'' boxes ask you to predict before reading on. Exercises are graded from $\star$
to $\star\star\star$; each has a hint and a solution that begins by explaining how one might find the idea.

Each chapter has a companion Python notebook that implements its algorithms from scratch, tests them against brute
force, and measures them against the bounds proved in the text, so that the reader can watch a $\Theta(n\log n)$
curve bend away from a $\Theta(n^2)$ one, or a randomised algorithm fail exactly as often as the theorem allows.
